% ============================================================
%  RiskTrail — Articulo JIC 2026
%  Compilar en Overleaf con pdfLaTeX (compilador por defecto)
% ============================================================
\documentclass[10pt,letterpaper,twocolumn]{article}

\usepackage[top=3cm,bottom=2.5cm,left=1.5cm,right=1.5cm,columnsep=1cm]{geometry}
\usepackage{newtxtext,newtxmath}
\usepackage[spanish,es-noshorthands]{babel}
\usepackage{amsmath}
\usepackage{booktabs}
\usepackage{array}
\usepackage{float}
\usepackage{graphicx}
\usepackage{caption}
\usepackage{microtype}
\usepackage{titlesec}

\captionsetup{font=footnotesize,labelfont=bf,justification=justified,singlelinecheck=false}

\setlength{\parindent}{0.5cm}
\setlength{\parskip}{0pt}

\titleformat{\section}{\fontsize{14}{16}\selectfont\bfseries}{\thesection.}{0.5em}{}
\titlespacing*{\section}{0pt}{8pt}{4pt}
\titleformat{\subsection}{\fontsize{10}{12}\selectfont\bfseries}{\thesubsection}{0.5em}{}
\titlespacing*{\subsection}{0pt}{6pt}{2pt}
\titleformat{\subsubsection}{\fontsize{10}{12}\selectfont\bfseries\itshape}{\thesubsubsection}{0.5em}{}
\titlespacing*{\subsubsection}{0pt}{4pt}{2pt}

% ============================================================
\begin{document}

% ---- BLOQUE DE TITULO (ancho completo) --------------------
\twocolumn[
  \begin{center}
    {\fontsize{18}{22}\selectfont\bfseries
    RiskTrail: sistema de evaluaci\'{o}n din\'{a}mica de riesgo para rutas
    de senderismo mediante modelos biom\'{e}c\'{a}nicos y datos
    meteorol\'{o}gicos en tiempo real\par}
    \vspace{12pt}
    {\fontsize{18}{22}\selectfont\bfseries
    RiskTrail: dynamic risk assessment system for hiking trails
    using biomechanical models and real-time meteorological data\par}
    \vspace{12pt}
    Roy [Apellido]$^{1}$,
    Kelvin He Wu$^{1}$,
    Irvin Ben\'{i}tez$^{1}$,
    [Nombre Asesor]$^{2*}$\par
    \vspace{6pt}
    {\fontsize{10}{12}\selectfont\itshape
    $^{1}$Facultad de Ingenier\'{i}a de Sistemas Computacionales,
    Universidad Tecnol\'{o}gica de Panam\'{a}, Ciudad de Panam\'{a},
    Panam\'{a}\\[2pt]
    $^{2}$[Unidad / Grupo de Investigaci\'{o}n],
    Universidad Tecnol\'{o}gica de Panam\'{a}, Ciudad de Panam\'{a},
    Panam\'{a}\par}
  \end{center}

  \vspace{8pt}
  \noindent\rule{\textwidth}{0.4pt}
  \vspace{6pt}

  % ---- RESUMEN -------------------------------------------
  \noindent{\fontsize{12}{14}\selectfont\textbf{Resumen}}\enspace
  El ecoturismo en ecosistemas neotropicales crece aceleradamente, pero
  las herramientas de planificaci\'{o}n de rutas disponibles emplean modelos
  est\'{a}ticos que omiten la interacci\'{o}n din\'{a}mica entre topograf\'{i}a,
  condiciones clim\'{a}ticas y gasto energ\'{e}tico humano. Este
  art\'{i}culo presenta \emph{RiskTrail}, un sistema web de evaluaci\'{o}n
  din\'{a}mica de riesgo sustentado en el \textbf{\'{I}ndice Din\'{a}mico
  de Riesgo (IDR)}, dise\~{n}ado para rutas de senderismo en Panam\'{a}.
  La metodolog\'{i}a combin\'{o} una revisi\'{o}n sistem\'{a}tica de
  literatura para fundamentar los modelos cient\'{i}ficos con el dise\~{n}o
  e implementaci\'{o}n de un prototipo funcional. El sistema integra
  archivos espaciales GPX, datos meteorol\'{o}gicos en tiempo real de la
  API Open-Meteo (temperatura, humedad relativa, precipitaci\'{o}n) y
  modelos biom\'{e}c\'{a}nicos validados: el Polinomio de Minetti para
  costo metab\'{o}lico, la funci\'{o}n de Tobler--Irmischer para velocidad
  predictiva y la ecuaci\'{o}n de Pandolf para esfuerzo con carga. La
  arquitectura comprende una capa de datos espaciales sobre
  PostgreSQL/PostGIS con pgRouting para enrutamiento topol\'{o}gico, un
  motor de c\'{a}lculo as\'{i}ncrono en FastAPI y una interfaz en Vue.js
  que renderiza un mapa de calor del riesgo. Los resultados del prototipo
  validan la viabilidad t\'{e}cnica de la arquitectura propuesta,
  clasificando segmentos en cinco niveles de riesgo alineados con la
  metodolog\'{i}a MIDE.\par
  \vspace{4pt}
  \noindent\textbf{Palabras clave}\enspace
  Biom\'{e}c\'{a}nica, ecoturismo, \'{i}ndice de riesgo, pgRouting,
  senderismo.\par
  \vspace{10pt}

  % ---- ABSTRACT ------------------------------------------
  \noindent{\fontsize{12}{14}\selectfont\textbf{Abstract}}\enspace
  Ecotourism in neotropical ecosystems is growing rapidly, yet current
  route planning tools rely on static models that neglect the dynamic
  interaction between topography, weather conditions, and human energy
  expenditure. This paper presents \emph{RiskTrail}, a dynamic risk
  assessment web system based on the \textbf{Dynamic Risk Index (DRI)},
  designed for hiking trails in Panama. The methodology combined a
  systematic literature review to ground the scientific models with the
  design and implementation of a functional prototype. The system
  integrates GPX spatial files, real-time meteorological data from the
  Open-Meteo API (temperature, relative humidity, precipitation), and
  validated biomechanical models: the Minetti Polynomial for metabolic
  cost, the Tobler--Irmischer hiking function for predictive walking
  speed, and the Pandolf equation for effort on loaded terrain. The
  architecture comprises a spatial data layer on PostgreSQL/PostGIS with
  pgRouting for topological routing, an asynchronous computation engine
  in FastAPI, and a Vue.js visualization interface rendering an
  interactive route risk heat map. Prototype results validate the
  technical feasibility of the proposed architecture, classifying route
  segments into five risk levels aligned with the MIDE methodology.\par
  \vspace{4pt}
  \noindent\textbf{Keywords}\enspace
  Biomechanics, dynamic risk index, ecotourism, hiking trails,
  pgRouting.\par
  \vspace{10pt}
  \noindent\rule{\textwidth}{0.4pt}
  \vspace{4pt}
  {\fontsize{8}{10}\selectfont\itshape
  $^{*}$\,Correo de contacto: [correo\_asesor@utp.ac.pa]\par}
  \vspace{12pt}
]

% ============================================================
%  CUERPO — 2 columnas
% ============================================================

\section{Introducci\'{o}n}

El senderismo se ha consolidado como una de las principales modalidades
del ecoturismo en Am\'{e}rica Latina, dinamizando econom\'{i}as locales y
promoviendo la conservaci\'{o}n de ecosistemas naturales. Panam\'{a}, por
su riqueza orogr\'{a}fica y biodiversidad neotropical, posee un alto
potencial para esta actividad~[14]: senderos como el Camino de Cruces y el
Parque Nacional Altos de Campana atraen miles de excursionistas anualmente.
Sin embargo, este crecimiento plantea un desaf\'{i}o fundamental: garantizar
la seguridad de los visitantes en entornos de alta variabilidad
clim\'{a}tica y topogr\'{a}fica.

Las plataformas actuales de planificaci\'{o}n ---como AllTrails o
Wikiloc--- presentan un enfoque est\'{a}tico: indican distancia, desnivel
y altitud, pero ignoran c\'{o}mo las condiciones ambientales afectan el
cuerpo humano durante el recorrido. Un sendero catalogado como «moderado»
en condiciones secas puede transformarse en un trayecto de alto riesgo tras
precipitaciones tropicales, cuando la combinaci\'{o}n de barro, pendientes
pronunciadas y humedad extrema eleva la probabilidad de resbalones, ca\'{i}das
y agotamiento prematuro. Estudios sobre tribolog\'{i}a y est\'{a}bilidad
postural confirman que la degradaci\'{o}n del coeficiente de fricci\'{o}n
en superficies h\'{u}medas es una causa primaria de ca\'{i}das durante la
locomoci\'{o}n~[4].

\subsection{Estado del arte}

La revisi\'{o}n de la literatura biom\'{e}c\'{a}nica identifica modelos de
alta precisi\'{o}n para estimar el gasto energ\'{e}tico en funci\'{o}n del
gradiente topogr\'{a}fico. La estimaci\'{o}n de tiempos de marcha se
fundament\'{o} hist\'{o}ricamente en la regla de Naismith~[8], que asume
una velocidad base lineal pero ignora la asimetr\'{i}a del descenso. El
Polinomio de Minetti~[1] super\'{o} esta limitaci\'{o}n al describir una
curva asim\'{e}trica del costo metab\'{o}lico que diferencia el esfuerzo de
ascenso del de descenso. Los modelos de velocidad de Tobler~[5] e
Irmischer-Clarke~[6] estiman tiempos de recorrido corrigiendo anomal\'{i}as
de los enfoques lineales, y trabajos recientes con telemetr\'{i}a GPS a gran
escala~[7] validan la variabilidad individual de la marcha seg\'{u}n el
perfil del excursionista.

El clima tropical de Panam\'{a} exacerba los riesgos: temperaturas
superiores a 32\,\textdegree C con humedad relativa mayor al 80\,\%
desencadenan deriva cardiovascular~[3], reduciendo la capacidad aer\'{o}bica
hasta un 40\,\% m\'{a}s r\'{a}pido que en condiciones temperadas~[2]. En el
\'{a}mbito de los Sistemas de Informaci\'{o}n Geogr\'{a}fica (SIG),
PostgreSQL con PostGIS~[9] y pgRouting~[10] permite recalcular costos en una
red topol\'{o}gica asimilando datos din\'{a}micos en tiempo real, enfoque
validado en sistemas de enrutamiento de emergencia~[11].

A pesar de la madurez individual de estos modelos, no se identific\'{o} una
plataforma que los integre en un \'{i}ndice unificado y din\'{a}mico
aplicado a entornos neotropicales. Este trabajo presenta \textbf{RiskTrail},
un sistema web que implementa el \textbf{\'{I}ndice Din\'{a}mico de Riesgo
(IDR)}, integrando dichos modelos biom\'{e}c\'{a}nicos, datos
meteorol\'{o}gicos en tiempo real y enrutamiento espacial adaptativo.

\section{Objetivos}

\textbf{Objetivo general:} Dise\~{n}ar e implementar un sistema web de
evaluaci\'{o}n din\'{a}mica de riesgo para rutas de senderismo que integre
modelos biom\'{e}c\'{a}nicos y datos meteorol\'{o}gicos en tiempo real,
aplicable a ecosistemas neotropicales paname\~{n}os.

\textbf{Objetivos espec\'{i}ficos:}
\begin{itemize}
  \setlength\itemsep{0pt}
  \item Fundamentar mediante revisi\'{o}n sistem\'{a}tica los modelos
        cient\'{i}ficos de costo metab\'{o}lico, velocidad de marcha y
        estr\'{e}s t\'{e}rmico aplicables al senderismo.
  \item Implementar los modelos de Minetti, Tobler--Irmischer y Pandolf
        para calcular costo metab\'{o}lico y velocidad predictiva por
        segmento de ruta.
  \item Integrar datos climatol\'{o}gicos en tiempo real mediante
        Open-Meteo para modelar el estr\'{e}s t\'{e}rmico (WBGT) y la
        fricci\'{o}n tribol\'{o}gica en descensos mojados.
  \item Construir un motor de enrutamiento topol\'{o}gico con pgRouting
        con costos asim\'{e}tricos en los arcos del grafo de red.
  \item Visualizar el IDR como mapa de calor interactivo en Vue.js con
        MapLibre GL JS.
\end{itemize}

\section{Materiales y M\'{e}todos}

\subsection{Revisi\'{o}n sistem\'{a}tica de literatura}

La identificaci\'{o}n de los modelos cient\'{i}ficos que sustentan el IDR
se realiz\'{o} mediante una b\'{u}squeda sistem\'{a}tica en \textbf{Google
Scholar}, \textbf{Semantic Scholar} y \textbf{Connected Papers}, empleando
t\'{e}rminos clave asociados a la biomec\'{a}nica de la marcha, el costo
metab\'{o}lico de transporte, las funciones de velocidad en senderismo y el
enrutamiento espacial din\'{a}mico. El cribado, la organizaci\'{o}n
tem\'{a}tica y la s\'{i}ntesis de los documentos recuperados se apoyaron en
\textbf{NotebookLM}, que permiti\'{o} extraer de forma estructurada los
hallazgos relevantes de cada fuente y descartar trabajos sin pertinencia
metodol\'{o}gica. La formalizaci\'{o}n matem\'{a}tica de los modelos
seleccionados ---Minetti, Tobler--Irmischer y Pandolf--- y su adaptaci\'{o}n
algor\'{i}tmica al \'{I}ndice Din\'{a}mico de Riesgo se asisti\'{o} con el
modelo de lenguaje \textbf{Claude (Opus)}, verificando en cada caso la
correspondencia entre las ecuaciones implementadas en el sistema y sus
fuentes cient\'{i}ficas originales.

\subsection{Arquitectura del sistema}

RiskTrail adopta una arquitectura desacoplada en tres capas. La
\textbf{capa de datos} emplea PostgreSQL~16 con PostGIS~3.4 y
pgRouting~3.6 para almacenamiento y an\'{a}lisis de geometr\'{i}as
tridimensionales en SRID~4326. La \textbf{capa l\'{o}gica} implementa
un microservicio as\'{i}ncrono en FastAPI (Python~3.12) con SQLAlchemy y
GeoAlchemy2, organizado en m\'{o}dulos de dominio independientes
(\emph{vel}, \emph{met}, \emph{rut}, \emph{cli}, \emph{rie}). La
\textbf{capa de presentaci\'{o}n} es una Aplicaci\'{o}n de P\'{a}gina
\'{U}nica en Vue~3 con TypeScript, Pinia, MapLibre GL~JS y Tailwind CSS.
La tabla~1 resume el \textit{stack} tecnol\'{o}gico.

\begin{table}[H]
\centering
\caption{\textbf{Tabla 1.} Stack tecnol\'{o}gico de RiskTrail}
\label{tab:stack}
\begin{tabular}{@{}ll@{}}
\toprule
\textbf{Capa} & \textbf{Tecnolog\'{i}a} \\
\midrule
Base de datos & PostgreSQL 16, PostGIS, pgRouting \\
Backend       & FastAPI, SQLAlchemy, GeoAlchemy2 \\
Frontend      & Vue 3, MapLibre GL JS, Pinia \\
Contenedores  & Docker, Docker Compose \\
Datos clima   & API Open-Meteo (c\'{o}digo abierto) \\
\bottomrule
\end{tabular}
\end{table}

\subsection{Modelos biom\'{e}c\'{a}nicos}

\subsubsection{Velocidad predictiva}

La velocidad de marcha en funci\'{o}n del gradiente~$i$ (decimal) se
estima con la funci\'{o}n gaussiana de Irmischer y Clarke~[6]:
\begin{equation}
  v(i) = \!\left[0.11 + e^{-\frac{(i + 0.0506)^{2}}{2\cdot 0.2043^{2}}}\right]\!\cdot 3.6
  \label{eq:irmischer}
\end{equation}
donde $v$ se expresa en km/h. A diferencia del modelo exponencial de
Tobler~[5], la formulaci\'{o}n gaussiana evita velocidades irreales en
gradientes extremos, condici\'{o}n habitual en senderos paname\~{n}os.

\subsubsection{Costo metab\'{o}lico de transporte}

El costo de transporte $C(i)$ en J/(kg$\cdot$m) se calcula con el
Polinomio de Minetti~[1], v\'{a}lido para $i\in[-0.45,\,{+}0.45]$:
\begin{equation}
  C(i) = 280.5i^{5} - 58.7i^{4} - 76.8i^{3} + 51.9i^{2} + 19.6i + 2.5
  \label{eq:minetti}
\end{equation}

El m\'{i}nimo de costo se alcanza en $i\approx{-}0.10$
($\approx$0.81\,J/(kg$\cdot$m)); en descensos abruptos ($i={-}0.45$)
asciende a 3.46\,J/(kg$\cdot$m) por contracciones exc\'{e}ntricas,
fuente primaria de lesiones articulares.

Para excursionistas con carga $L$ (kg) se aplica la ecuaci\'{o}n de
Pandolf~[15]:
\begin{equation}
  M = 1.5W + 2.0(W\!+\!L)\!\left(\!\frac{L}{W}\!\right)^{\!2}
      + \eta(W\!+\!L)(1.5V^{2} + 0.35VG)
  \label{eq:pandolf}
\end{equation}
donde $W$ es el peso corporal (kg), $V$ la velocidad (m/s), $G$ el
gradiente (\%) y $\eta$ el coeficiente de terreno (1.0 asfalto;
1.2 tierra; 2.1 barro).

\subsection{Modelado del estr\'{e}s clim\'{a}tico}

El estr\'{e}s t\'{e}rmico se cuantifica con el \'{I}ndice de Globo de
Bulbo H\'{u}medo (WBGT):
\begin{equation}
  \mathrm{WBGT} = 0.7T_{w} + 0.2T_{g} + 0.1T_{a}
  \label{eq:wbgt}
\end{equation}
donde $T_{w}$, $T_{g}$ y $T_{a}$ son las temperaturas de bulbo
h\'{u}medo, globo y ambiente. Cuando el excursionista supera 20 minutos
de marcha bajo un WBGT mayor al umbral cr\'{i}tico, el algoritmo aplica
una penalizaci\'{o}n exponencial al costo temporal, modelando la deriva
cardiovascular documentada por Rowell~[3].

El riesgo tribol\'{o}gico en descensos con pendiente menor a
$-10{^\circ}$ y precipitaci\'{o}n reciente activa un multiplicador
$\mu_{\mathrm{eff}} = \mu_{\mathrm{base}}\cdot f(<1)$ que incrementa el
costo de los arcos de descenso en pgRouting, reflejando la reducci\'{o}n
del coeficiente de fricci\'{o}n en suelo h\'{u}medo~[4].

\subsection{Procesamiento espacial}

Los archivos GPX se transforman en \textit{linestrings} 3D en la base de
datos. La funci\'{o}n \texttt{ST\_3DDistance} de PostGIS extrae la
distancia planim\'{e}trica y el diferencial de elevaci\'{o}n entre
v\'{e}rtices para calcular gradientes exactos. pgRouting construye un
grafo topol\'{o}gico con costos asim\'{e}tricos: el costo directo emplea
$C({+}i)$ de la ecuaci\'{o}n~(\ref{eq:minetti}) y el costo inverso
$C({-}i)$, reflejando la asimetr\'{i}a fisiol\'{o}gica ascenso--descenso.
Los datos meteorol\'{o}gicos se recuperan de forma as\'{i}ncrona desde
Open-Meteo en el centroide de cada segmento.

\subsection{Clasificaci\'{o}n IDR y escala MIDE}

El IDR integra cin\'{e}matica (ec.~\ref{eq:irmischer}),
metabolismo (ecs.~\ref{eq:minetti}--\ref{eq:pandolf}) y entorno
(ec.~\ref{eq:wbgt}) en una puntuaci\'{o}n por segmento, mapeada a cinco
niveles alineados con la metodolog\'{i}a MIDE~[13] (tabla~2).

\begin{table}[H]
\centering
\caption{\textbf{Tabla 2.} Escala del \'{I}ndice Din\'{a}mico de Riesgo}
\label{tab:idr}
\begin{tabular}{@{}clp{2.5cm}@{}}
\toprule
\textbf{Nivel} & \textbf{Condici\'{o}n} & \textbf{Color / Acci\'{o}n} \\
\midrule
1 & Sin riesgos relevantes     & Verde / Sin restricci\'{o}n \\
2 & Fatiga moderada esperada   & Amarillo / Recomendaci\'{o}n \\
3 & Estr\'{e}s t\'{e}rmico elevado     & Naranja / Alerta \\
4 & Riesgo tribol\'{o}gico alto & Rojo / Advertencia \\
5 & Peligro extremo combinado  & Morado / Restricci\'{o}n \\
\bottomrule
\end{tabular}
\end{table}

\section{Resultados}

El prototipo de RiskTrail fue validado con archivos GPX de rutas
representativas de parques nacionales paname\~{n}os sobre el sistema
desplegado en contenedores Docker. La tabla~3 resume el desempe\~{n}o
comparativo de los tres modelos de velocidad implementados.

\begin{table}[H]
\centering
\caption{\textbf{Tabla 3.} Comparaci\'{o}n de modelos de velocidad}
\label{tab:modelos}
\begin{tabular}{@{}lccc@{}}
\toprule
\textbf{Modelo} & \textbf{Grad. \'{o}ptimo} & \textbf{$v$ m\'{a}x.} & \textbf{Extremos} \\
\midrule
Naismith  & 0\,\%    & 5.0\,km/h & No modelado \\
Tobler    & $-$5\,\% & 6.0\,km/h & Irreal \\
Irmischer & $-$5\,\% & 5.9\,km/h & V\'{a}lido \\
\bottomrule
\end{tabular}
\end{table}

La API REST entreg\'{o} respuestas de an\'{a}lisis IDR en tiempos menores
a 1.2\,s para rutas de hasta 500 segmentos, validando la eficacia de la
arquitectura as\'{i}ncrona. El modelo de Minetti confirm\'{o} la
asimetr\'{i}a costo--pendiente: un segmento con gradiente ${+}0.30$
registr\'{o} 10.8\,J/(kg$\cdot$m), mientras el equivalente en descenso
(${-}0.30$) arroj\'{o} 3.1\,J/(kg$\cdot$m), indicando que los descensos
pronunciados generan riesgo de lesi\'{o}n por contracci\'{o}n
exc\'{e}ntrica y no por fatiga aer\'{o}bica acumulada.

La interfaz Vue.js renderiz\'{o} correctamente el mapa de calor del IDR,
diferenciando mediante interpolaci\'{o}n de color en MapLibre GL~JS los
segmentos seguros (verde, nivel~1) de los tramos con riesgo combinado de
estr\'{e}s t\'{e}rmico y terreno h\'{u}medo en descenso (rojo--morado,
niveles~4--5). La escala MIDE qued\'{o} integrada como m\'{e}trica de
salida normalizada visible al usuario en la barra lateral.

\section{Discusi\'{o}n}

Los resultados validan que la integraci\'{o}n de modelos
biom\'{e}c\'{a}nicos con datos espaciales y meteorol\'{o}gicos en tiempo
real supera el enfoque de las plataformas est\'{a}ticas existentes.
A diferencia de AllTrails o Wikiloc, RiskTrail clasifica el riesgo por
segmento individual, permitiendo identificar los tramos espec\'{i}ficos
que requieren precauci\'{o}n.

La elecci\'{o}n del modelo de Irmischer sobre Tobler puro resulta
justificada en gradientes extremos: la funci\'{o}n gaussiana evita
velocidades predichas irreales en pendientes superiores a
$35{^\circ}$ o inferiores a $-35{^\circ}$, condici\'{o}n frecuente en
senderos paname\~{n}os. Campbell et al.~[7] confirman que la variabilidad
individual en la velocidad depende significativamente del gradiente y el
perfil del excursionista, justificando futuras extensiones del modelo con
datos personalizados.

La incorporaci\'{o}n del WBGT como disparador de la deriva cardiovascular
es especialmente relevante para el clima tropical paname\~{n}o. Maughan
et al.~[2] confirman que en entornos con humedad absoluta alta el
agotamiento se instala hasta un 40\,\% m\'{a}s r\'{a}pido que en
condiciones temperadas, factor ignorado por las herramientas actuales.

Las limitaciones del prototipo incluyen la ausencia de validaci\'{o}n en
campo con excursionistas reales y la dependencia de la resoluci\'{o}n del
Modelo Digital de Elevaci\'{o}n (DEM) para la precisi\'{o}n de los
gradientes calculados.

\section{Conclusiones}

RiskTrail demuestra la viabilidad t\'{e}cnica de un sistema de
evaluaci\'{o}n din\'{a}mica de riesgo que integra modelos
biom\'{e}c\'{a}nicos validados, datos meteorol\'{o}gicos en tiempo real y
enrutamiento espacial adaptativo para rutas de senderismo. El IDR
clasifica correctamente los segmentos de mayor riesgo en funci\'{o}n de
la topograf\'{i}a y las condiciones clim\'{a}ticas, superando el enfoque
est\'{a}tico de las herramientas existentes.

El sistema contribuye a la seguridad de excursionistas en ecosistemas
neotropicales paname\~{n}os donde la variabilidad clim\'{a}tica es extrema.
La arquitectura modular es escalable: permite incorporar nuevos perfiles de
usuario, sensores IoT y la metodolog\'{i}a de capacidad de carga
tur\'{i}stica de Cifuentes~[12] como capa adicional de gesti\'{o}n.

El trabajo futuro incluye la validaci\'{o}n en campo en el Parque Nacional
Camino de Cruces, la integraci\'{o}n de telemetr\'{i}a GPS en tiempo real
y el desarrollo de un m\'{o}dulo de recomendaci\'{o}n de rutas
\'{o}ptimas ponderadas por perfil del excursionista mediante pgRouting.

\section*{AGRADECIMIENTOS}

A la Universidad Tecnol\'{o}gica de Panam\'{a} por el apoyo institucional
durante el desarrollo de esta investigaci\'{o}n. A los equipos de
PostGIS, pgRouting y Open-Meteo por mantener herramientas de c\'{o}digo
abierto que fundamentaron la implementaci\'{o}n t\'{e}cnica del sistema.

% ============================================================
%  REFERENCIAS (formato IEEE, 9pt)
% ============================================================
\section*{REFERENCIAS}
{\fontsize{9}{11}\selectfont
\begin{enumerate}
  \setlength\itemsep{2pt}
  \setlength\parskip{0pt}

  \item A. E. Minetti, C. Moia, G. S. Roi, D. Susta y G. Ferretti,
        ``Energy cost of walking and running at extreme uphill and downhill
        slopes,'' \textit{J. Appl. Physiol.}, vol.~93, no.~3,
        pp.~1039--1046, 2002.

  \item R. J. Maughan, J. E. Wingo y E. F. Coyle, ``Delineating the
        impacts of air temperature and humidity for endurance exercise,''
        \textit{Exp. Physiol.}, vol.~107, no.~5, pp.~386--395, 2022.

  \item L. B. Rowell, ``Cardiovascular drift during prolonged exercise:
        new perspectives,'' \textit{Med. Sci. Sports Exerc.}, vol.~15,
        no.~4, pp.~367--379, 1983.

  \item T. E. Lockhart \textit{et al.}, ``Slips and falls,'' Arizona State
        University, Ira A. Fulton Schools of Engineering, 2016.

  \item W. Tobler, ``Three presentations on geographical analysis and
        modeling,'' National Center for Geographic Information and
        Analysis, Informe T\'{e}cnico 93-1, 1993.

  \item I. J. Irmischer y K. C. Clarke, ``Hiking with Tobler: tracking
        movement and calibrating a cost function for personalized 3D
        accessibility,'' \textit{Findings}, 2018.

  \item M. J. Campbell, P. E. Dennison y M. P. Thompson, ``Predicting the
        variability in pedestrian travel rates and times using crowdsourced
        GPS data,'' \textit{Comput. Environ. Urban Syst.}, vol.~97,
        p.~101866, 2022.

  \item W. Naismith, ``Naismith's rule,'' \textit{Scottish Mountaineering
        Club Journal}, vol.~2, no.~3, p.~136, 1892.

  \item The PostGIS Development Team, ``PostGIS 3.4 Reference Manual,''
        2024. [En l\'{i}nea]. Disponible: https://postgis.net/docs/

  \item pgRouting Development Team, ``pgRouting 3.6 Documentation,''
        2024. [En l\'{i}nea]. Disponible: https://docs.pgrouting.org/

  \item S. Choosumrong, R. Raghavan, V. Pundt y N. Intajag,
        ``Implementation of dynamic routing as a web service for
        emergency routing decision planning,'' \textit{Int. Arch.
        Photogramm. Remote Sens. Spatial Inf. Sci.}, vol.~XL-2,
        pp.~59--63, 2014.

  \item M. Cifuentes, ``Determinaci\'{o}n de capacidad de carga
        tur\'{i}stica en \'{a}reas protegidas,'' CATIE, Turrialba,
        Informe T\'{e}cnico 194, 1992.

  \item Monta\~{n}a Segura, ``Manual MIDE: M\'{e}todo para
        informaci\'{o}n de excursiones,'' Federaci\'{o}n Espa\~{n}ola de
        Deportes de Monta\~{n}a y Escalada, s.f.

  \item Geoversity, ``GIS in the jungle: Experiential Environmental
        Education (EEE) in Panama,'' 2021.

  \item K. B. Pandolf, B. Givoni y R. F. Goldman, ``Predicting energy
        expenditure with loads while standing or walking very slowly,''
        \textit{J. Appl. Physiol.}, vol.~43, no.~4, pp.~577--581, 1977.

\end{enumerate}
}

\end{document}
